\subsection{Validation, timings and integration scope}

All 14 focused tests pass in 1.753 seconds; the full maintained suite passes
all 1,219 tests in 254.202 seconds. The final source audit takes 19.307
seconds. The audit and full suite partly overlap, so those elapsed times
are correctness-run observations rather than performance comparisons. The
final audit and both timing runs pin the same 501 source files. Existing
source-bound disc, annulus, planar, span and weighted-orbit verifiers remain
byte-identical to the baseline revision \code{3c0817c50}.

The performance endpoint starts from a diagram and includes finite exterior
construction, primitive cocycle extraction, minimum-span optimization, the
selected search and independent source-bound replay of its selected surface.
The old arm searches all distinct root extrema, using the number of global
vertices as its root allowance, and retains the highest Euler value found.
The new arm solves the whole face and checks its arithmetic dual. This is a
complete \emph{candidate-construction} pipeline, not complete knot
recognition. Earlier native successes are deliberately not allowed to skip
the experiment, and no exact knot-recognition fallback is timed here.

The first five-round run contains 120 measured calls and 24 warm-ups and
takes 24.281 seconds. Variation in the original-disc and Kinoshita--Terasaka
cases motivates the retained nine-round repeat: 216 measured calls and 24
warm-ups in 40.195 seconds. Each round randomizes roots, roots A/A, exact and
exact A/A calls. No tests, audits, builds or commits overlap either timing
run. Every measured call completes, and all per-arm Euler values, piece
counts, operation counts and certificate hashes are stable within and
between both runs. Primary-arm median times and median paired ratios from
the final run follow; paired ratios need not equal quotients of the displayed
medians.

\begin{center}
\small
\begin{tabular}{lrrrrr}
\toprule
Input & Roots (ms) & Exact (ms) & Roots/exact & Roots A/A & Exact A/A \\
\midrule
Trefoil & 38.209 & 37.957 & 1.021 & 0.940 & 0.976 \\
Original disc & 48.704 & 71.274 & 0.690 & 1.004 & 1.003 \\
Tied disc & 115.241 & 115.147 & 0.989 & 0.986 & 0.990 \\
Random 19 & 228.711 & 218.924 & 1.063 & 1.001 & 1.002 \\
Random 158 & 333.819 & 259.215 & 1.283 & 0.996 & 0.994 \\
Kinoshita--Terasaka & 202.822 & 227.087 & 0.825 & 1.027 & 0.995 \\
\bottomrule
\end{tabular}

\end{center}

Random case 158 improves by paired factors 1.272 and 1.283 in the two runs;
random case 19 improves by 1.054 and 1.062. The original-disc case is slower,
with ratios 0.677 and 0.690, and Kinoshita--Terasaka is slower with ratios
0.835 and 0.825. Those final ratios correspond to about 45 and 21 percent
additional time. The tied-disc case is near parity. The trefoil's small
apparent improvement is not evidence of a reliable gain: its final roots
A/A ratio is 0.940. Absolute Kinoshita--Terasaka times also vary substantially
between runs. The data support mixed costs, not a uniform speedup.

The new callable \code{maximize\_cocycle\_face\_euler} API is consequently
available for explicit candidate construction and research; it is not added
to the default native schedule. All existing candidate and portfolio defaults
remain unchanged. Its theoretical gain is an exact optimum over a face described by $8N$
inequalities, together with concrete certified limitations of that face. It supplies neither all normal-surface candidates
nor a complete planar-capping search, a hierarchy algorithm, or a general
subexponential or quasipolynomial recognition bound.

Reproduce with \code{normal\_orbit\_research.euler} in \code{audit} mode or
\code{benchmark --rounds 9} mode. The indexed models and arithmetic duals
are retained in the audit JSON for direct independent replay. The incoming
unknot archive inventory and the checked \code{openai/math} head
\code{fd4aeeb2e} remain unchanged. The new reduction is derived from the
maintained height-cell identity and matching-face description above.
